% 第3章第7节：本章小结。

\section{本章小结}
\label{sec:pde-summary}

本章开头问了四个问题：哪些数据确定一个场，粗糙数据下“满足方程”是什么意思，三种演化机制各自做了什么，网格上的近似何时可靠。下面按这四个问题收束。

\paragraph{四个代表方程的对照}
表~\ref{tab:pde-four-equations}把四个方程放在一起。Poisson方程描述平衡，在整个边界上给数据，解属于$H^1$类空间，Lax--Milgram定理与能量估计给出存在、唯一与连续依赖；边界条件一变，解空间、相容条件和唯一性随之改变。热方程描述耗散：它需要初值和整个边界上的数据，正弦模式按$e^{-k^2\pi^2t}$衰减，高频先消失，解在正时刻立即变光滑，能量单调减少，最大值原理成立，反向求解不适定。输运方程描述搬运：沿特征线它化为常微分方程，只在流入边界给数据，间断被原样搬走；守恒形式还要计入流动的压缩，非线性时特征线可能相交并产生激波。波动方程描述传播：需要位移与速度两个初值，信号以有限速度$c$传播，能量守恒，粗糙的初值永远不会被磨平，反向求解与正向一样适定。

\begin{table*}[!htb]
  \centering
  \begin{tabularx}{\linewidth}{@{}l>{\RaggedRight\arraybackslash}X>{\RaggedRight\arraybackslash}X>{\RaggedRight\arraybackslash}X>{\RaggedRight\arraybackslash}X@{}}
    \toprule
    方程 & 所需数据 & 解的自然空间 & 关键估计或性质 & 信息传递 \\
    \midrule
    Poisson $-\Delta u=f$ & 源项，整个边界 & $H_0^1$或$H^1$ & 强制性与能量估计 & 平衡，无时间 \\
    热方程 $u_t=\kappa\Delta u+f$ & 初值，整个边界 & 时间上取值于$H_0^1$ & 能量耗散，最大值原理 & 无穷速度，磨光 \\
    输运 $u_t+\bm b\cdot\nabla u=0$ & 初值，流入边界 & 沿特征线的$C^1$或弱解 & 沿特征线保值 & 沿特征线，不磨光 \\
    波动 $u_{tt}=c^2\Delta u$ & 初位移与初速度，整个边界 & $(u,u_t)\in H_0^1\times L^2$ & 能量守恒 & 有限速度$c$，不磨光 \\
    \bottomrule
  \end{tabularx}
  \caption{本章四个代表方程的对照（线性情形，有界区域或全直线）。}
  \label{tab:pde-four-equations}
\end{table*}

\paragraph{连续场、弱方程与有限维近似}
贯穿全章的工具是弱形式：用测试函数检验方程，把导数转移到测试函数上。它让源项可以只是可积函数甚至一维中的点源，让边界条件分成写进解空间的本质条件和藏在边界积分里的自然条件，也让能量估计成为证明唯一性与连续依赖的统一方法。同一个弱形式又直接给出数值方法：把解空间和检验空间同时换成有限维子空间，就是Galerkin与有限元方法；Galerkin正交性和Céa引理把数值误差归结为子空间的逼近能力。有限差分从另一端出发，用截断误差衡量一致性，用离散最大值原理证明稳定性，二者合起来给出收敛。演化问题经空间离散后成为刚性的常微分方程组，显式格式受$\Delta t\leq h^2/(2\kappa)$或CFL条件$c\Delta t\leq h$的限制，隐式格式解除了稳定性限制，但需要按精度平衡空间与时间误差。

\paragraph{几种不同的命题}
本章反复区分了几类结论，它们不能互相代替。写出弱形式或Euler--Lagrange方程，只说明“若有解，它满足什么”，不等于证明了解存在；存在与唯一需要Lax--Milgram定理这样的论证。形式级数只说明模式怎样组合，在什么空间中收敛、能否逐项求导，取决于初值的光滑程度。能量等式说明某个量守恒或衰减，它给出唯一性与连续依赖，但不直接给出逐点的光滑性。截断误差小只说明精确解几乎满足网格方程；离散解收敛到连续解，还需要与网格无关的稳定估计。残差小也只有经过相应的稳定估计，才能解释为误差小。

\paragraph{与后续章节的联系}
本章始终处理确定性的场。第~\ref{chap:continuous-time-stochastic-processes}章将看到，随机状态的概率密度怎样随时间演化：只有漂移时，密度满足本章的连续性方程；加入Brown运动的随机扰动后，密度满足带扩散项的方程，热方程就是其中最简单的情形。届时需要的确定性方程、弱解与半群语言都已在本章建立，随机过程一侧只需补充生成元与转移概率之间的联系。
